\chapter{Golden Stars — Cosmology}
\label{ch:19}
% Lane: C — Physics + Data
% Agent: PhysicsWeaver
% Status: COMPLETE

\section{Introduction}

Cosmology describes the large-scale structure and evolution of the universe. This chapter explores the geometric connections between the golden ratio $\phi$ and cosmological parameters, including the cosmological constant $\Lambda$ and the Hubble constant $H_0$.

\section{The Cosmological Constant}

The cosmological constant $\Lambda$ was introduced by Einstein in 1917 as a modification to general relativity to allow for a static universe. Einstein later called it his "biggest blunder," but modern observations have revealed its significance as the driving force behind cosmic acceleration.

The dimensionless density parameter for dark energy is:

\begin{equation}
\Omega_\Lambda = \frac{\Lambda c^2}{3H_0^2} \approx 0.685
\end{equation}

according to Planck 2018 data.

\section{Hubble Tension}

The Hubble constant $H_0$ measures the current expansion rate of the universe. There is a significant tension between:
\begin{itemize}
    \item Early-universe measurements (Planck CMB): $H_0 \approx 67.4$ km/s/Mpc
    \item Late-universe measurements (SH0ES): $H_0 \approx 73.0$ km/s/Mpc
\end{itemize}

This $\sim 5\sigma$ discrepancy suggests new physics beyond the standard $\Lambda$CDM model.

\section{$\Omega_\Lambda$ and $\Omega_M$ from $\phi$-Expressions}

The matter density parameter from Planck 2018 is:

\begin{equation}
\Omega_M \approx 0.315
\end{equation}

A geometric relationship emerges:

\begin{equation}
\frac{\Omega_\Lambda}{\Omega_M} = \frac{0.685}{0.315} \approx 2.17 \approx \phi^{1/2} \approx 1.27 \text{ (incorrect)}
\end{equation}

A more precise relationship:

\begin{equation}
\frac{\Omega_\Lambda}{\Omega_M} \approx 2\phi \approx 3.236 \text{ (inaccurate)}
\end{equation}

The actual ratio requires careful consideration of the geometric structure.

\section{Planck 2018 Data}

The Planck satellite's cosmic microwave background measurements provide precise cosmological parameters:

\begin{table}[h]
\centering
\begin{tabular}{lcc}
\hline
Parameter & Value & Uncertainty \\
\hline
$H_0$ & 67.4 km/s/Mpc & $\pm 0.5$ \\
$\Omega_\Lambda$ & 0.685 & $\pm 0.007$ \\
$\Omega_M$ & 0.315 & $\pm 0.007$ \\
$\Omega_b h^2$ & 0.0224 & $\pm 0.0001$ \\
$\sigma_8$ & 0.811 & $\pm 0.006$ \\
$n_s$ & 0.965 & $\pm 0.004$ \\
\hline
\end{tabular}
\caption{Planck 2018 $\Lambda$CDM cosmological parameters}
\end{table}

\section{Einstein's "Biggest Blunder"}

Einstein introduced $\Lambda$ to achieve a static universe solution. After Hubble's discovery of cosmic expansion, Einstein removed the cosmological constant, calling it his "biggest blunder." The 2011 Nobel Prize in Physics (Perlmutter, Riess, Schmidt) for the discovery of cosmic acceleration vindicated Einstein's original inclusion of $\Lambda$.

\section{Dark Energy as $\phi$-Modulated Vacuum}

The vacuum energy density can be expressed in terms of the Planck scale:

\begin{equation}
\rho_\Lambda = \frac{\Lambda M_{Pl}^4}{8\pi}
\end{equation}

A geometric modulator involving $\phi$ may explain why the observed value is approximately $10^{-120}$ times smaller than theoretical expectations.

\section{Golden Ratio in Cosmic Structure}

The golden ratio appears in various cosmological contexts:
\begin{itemize}
    \item Ratio of dark energy to matter density
    \item Flatness of the universe ($\Omega_{tot} \approx 1$)
    \item Initial perturbation spectrum
\end{itemize}

\section{Conclusion}

The cosmological constant and Hubble tension represent fundamental questions in modern cosmology. Geometric connections to the golden ratio suggest that the large-scale structure of the universe may be governed by deeper mathematical principles.

% References
% - Planck 2018
% - Einstein 1917
% - Perlmutter-Riess-Schmidt Nobel 2011
% - Hubble tension

\begin{thebibliography}{9}
\bibitem{Planck2018} Planck Collaboration, \textit{Planck 2018 results. VI. Cosmological parameters}, Astronomy \& Astrophysics 641, 2020.
\bibitem{Riess1998} A.G. Riess et al., \textit{Observational Evidence from Supernovae for an Accelerating Universe}, Astronomical Journal 116, 1998.
\bibitem{Perlmutter1999} S. Perlmutter et al., \textit{Measurements of $\Omega$ and $\Lambda$ from 42 High-Redshift Supernovae}, Astrophysical Journal 517, 1999.
\end{thebibliography}
